\subsection{The cuspidal criterion and its arithmetic inputs}

We retain the whole arithmetic Galois quotient throughout this subsection.

The manuscript question addressed is cote~149, scan~36: can the covering condition force the equality of the two points?

\begin{claim}{A.CR1}{Proved}{The same-prime quotient lemma}

Let $F$ be profinite, $K\triangleleft F$ closed, $r$ prime, and
$N=\ker(K\to K^{(r)})$. For every open $I\le F$ containing $N$, the
continuous quotient map induces a natural isomorphism
$I^{(r)}\simeq(I/N)^{(r)}$. No restriction on $[F:I]$ is imposed.

\end{claim}

\begin{proof}

Let F be a profinite group, K a closed normal subgroup, r a prime, and N the kernel of the maximal pro-r quotient \(K{\to}K^{(r)}\). Because N is characteristic in K, it is normal in F. Let \(I{\le}F\) be open and contain N. We prove that the continuous epimorphism \(I{\to}I/N\) induces an isomorphism \(I^{(r)}\to(I/N)^{(r)}\).

Take an arbitrary continuous quotient I/J which is a finite r-group, where J is open and normal in I. Put \(I_0=\operatorname{core}_F(I)\); it is open normal in F, contained in I, and contains N. Put \(J_0=J{\cap}I_0\). Then \(J_0\) is open normal in \(I_0\) and \(I_0/J_0\hookrightarrow I/J\), so \(I_0/J_0\) is an r-group. Take \(M=\operatorname{core}_F(J_0)\). There are only finitely many F-conjugates of \(J_0\), because \(J_0\) is open. Since \(I_0\) is normal in F, \(I_0/M\) embeds in the product of the quotients \(I_0/{}^fJ_0\), each isomorphic to \(I_0/J_0\). Thus \(I_0/M\) is an r-group.

The image of K in \(F/I_0\) is an r-group: \(N{\subset}I_0\), so this image is a finite quotient of K/N, a pro-r group. The image of K in F/M is an extension of that image by a subgroup of \(I_0/M\). It too is an r-group. By the universal defining property of N, N is killed by \(K{\to}F/M\). Therefore \(N\subset M\subset J_0\subset J\).

Every finite r-quotient of I consequently kills N. This is exactly the asserted isomorphism of maximal pro-r quotients, with its natural map, not merely an abstract isomorphism. No finite-generation assumption, normality of I, or r-power-index assumption on I was used.


\end{proof}

\begin{claim}{A.CR2}{Proved}{Compactification and Frobenius are retained at the cusp prime}

Let $C/k$ be smooth proper geometrically connected of genus at least two,
$k$ finite of characteristic $p$, $P\in C(k)$, and $r\ne p$.
Every open geometric subgroup of the maximal cuspidally pro-$r$
quotient of $\pi_1(C-P)^{(p')}$, after its cusp inertia is killed and
its maximal pro-$r$ quotient is taken, gives the full geometric
pro-$r$ fundamental group of its actual smooth compactification.
The comparison is compatible with every arithmetic action preserving
the cover. Thus a cusp-preserving isomorphism over $G_k$ preserves
the point counts of every corresponding open cover over every
finite constant extension.

\end{claim}

\begin{proof}

Take \(F=\pi_1((C-P)_{\bar k})^{(p')}\), \(K=\ker(F\to\Delta_C)\), \(r{\ne}p\). An open I/N in F/N corresponds to an actual connected finite étale cover \(D^\circ\to(C-P)_{\bar k}\) whose Galois closure has prime-to-p order. The full pro-prime-to-p topology on its geometric group is induced from F: for an open subgroup with a prime-to-p Galois core, the normal-core construction again gives a prime-to-p group, since extensions and subgroups of products of prime-to-p groups have prime-to-p order. Thus I is the full prime-to-p geometric group of \(D^{\circ}\).

Let D be its smooth compactification. Let \(T{\trianglelefteq}I\) be the closed normal subgroup generated by the inertia groups of the points in \(D-D^{\circ}\). Then \(I/T=\pi_1(D)^{(p')}\). The image of T in I/N is precisely the closed normal subgroup generated by the images of those inertia groups. Taking a maximal pro-r quotient commutes with quotient by a closed normal subgroup, by the universal property of continuous maps to finite r-groups. Therefore

\[
 ((I/N)/\overline{TN}/N)^{(r)}
 \simeq (I/T)^{(r)}
 =\pi_1(D)^{(r)}.
\]

This isomorphism is induced by the actual punctured-cover and compactification maps. Every arithmetic conjugation preserving the cover preserves its puncture inertia system and all kernels in this construction. It is therefore equivariant for the outer action of \(G_L\) when the cover descends to a finite field L. Consequently it identifies \(H^1(D_{\bar L},\mathbf Z_r)\) with its actual Frobenius action, not only its rank or characteristic polynomial modulo r.

If two such arithmetic relative groups are isomorphic over \(G_L\) and their cusp decomposition groups correspond, corresponding covers have isomorphic compactified r-adic \(H^1\) representations. Hence they have equal integral Weil polynomials and equal compactified point counts over each finite extension of L. The missing points are indexed by the cusp double cosets, and their degrees are the indices of the projected decomposition groups in \(G_L\). Thus these degrees also agree. Subtracting the missing points gives equal point counts for the open covers over every constant extension.

The point-count comparison uses the trace formula and the integral
Frobenius polynomial for a smooth proper curve over a finite field:
these are \cite[formula (1.5.1) and Theorem (1.6)]{WeilCurves}.
Their hypotheses apply to both smooth compactifications and to every
finite extension of their common constant field. Equality of their
\(r\)-adic characteristic polynomials is equality of integer polynomials,
since \(\mathbf Z[T]\hookrightarrow\mathbf Z_r[T]\).


\end{proof}

\begin{claim}{A.CR3}{Known, with derivation}{The exact uniform exceptional-set input}

For a proper hyperbolic curve $C/k$ over a finite field of
characteristic $p$, and $\Sigma=\mathrm{Primes}-\{p\}$, there is a
finite set $E_C\subset C^{\cl}$ such that outside its inverse image
a decomposition subgroup has a unique geometric point label, even
up to commensurability, over every finite constant extension at once.
The numbering of Definitions~1.3 and 1.5, Remark~1.2, and
Proposition~1.8(vi) is the same in the attached arXiv text and the
journal version of \cite{ST2009}.

\end{claim}

\begin{proof}

Let k be finite of characteristic p, X/k a smooth proper geometrically connected curve, \(U=X-S\) hyperbolic, and \({\Sigma}\) a set of primes containing at least one prime different from p.~Let \({\Pi}_U\) be the arithmetic group retaining the whole \(G_k\) and the maximal geometric pro-\({\Sigma}\) quotient of the tame group. Its universal pro-cover compactification is denoted by \(\widetilde X_U\).

The attached ST2009 Definition 1.3(iii) calls a point a of a set map f exceptional when \(\#f^{-1}(f(a))>1\). Definition 1.5 applies this to the map \(\overline D\) from closed points upstairs to commensurability classes of decomposition subgroups. Definition 1.10 puts \(E_U\) equal to its image in \(X^{\rm cl}\).

Proposition 1.8(iii) states injectivity of \(\overline D\) over a fixed geometric point; (iv) gives finite fibres over one closed point and injectivity over one rational point; (v) says the exceptional set is \({\Pi}_U\)-stable; under the additional hypothesis that \({\Sigma}\) is cofinite, (vi) states that \(E_{\widetilde U}/\Pi_U\) is finite. This quotient is \(E_U\). The same numbering and hypotheses occur in the opened journal PDF.

Remark 1.2 says that the commensurability maps are unchanged upon passing to any open subgroup. More explicitly, if H is open in \({\Pi}_U\) and \(D_1,D_2\) are closed subgroups, then \(D_1\) and \(D_2\) are commensurate if and only if \(D_1{\cap}H\) and \(D_2{\cap}H\) are commensurate. One direction is immediate. For the other, \(D_i{\cap}H\) has finite index in \(D_i\); the common intersection has finite index in both. Universal pro-cover points are the same on replacing U by its H-cover. In particular, when H is the preimage of \(G_L\) for a finite constant extension L/k, the exceptional points upstairs are exactly the same points and the exceptional image on \(X_L\) is the finite inverse image of \(E_U\).

For \({\Sigma}=\mathrm{Primes}-\{p\}\), all hypotheses are met when C is proper of genus at least two: U=C is hyperbolic and \({\Sigma}\) is nonempty and cofinite. Consequently there is one finite set \(E_C\subset C^{\rm cl}\) whose geometric inverse image excludes all exceptional labels over all finite extensions. If R is a closed point outside \(E_C\), a decomposition subgroup at R cannot be commensurate, even after conjugation, with a subgroup at another geometric point. Conjugation moves the lift and preserves the underlying closed-point label. This is exactly the use in AS-AG.40 Step 3.

\textbf{Check of their proof and its hypotheses.}

Their proof first replaces U by an open-subgroup cover whose compactification has genus \textgreater1 and is nonhyperelliptic. Commensurability invariance just proved makes finiteness upstairs descend under this finite cover. For a nonhyperelliptic compactification X, the difference morphism \(\delta:X\times X\to J_X\), \((P,Q)\mapsto[P-Q]\), is injective on geometric points away from the diagonal. Indeed, equality of two differences gives \(P+Q'\sim P'+Q\); distinct linearly equivalent effective degree-two divisors would supply a degree-two map to the projective line, contrary to nonhyperellipticity. Equality of the divisors then gives the equal ordered pair, unless both differences vanish.

The image W of \({\delta}\) contains no translate of a positive-dimensional abelian subvariety. A two-dimensional such translate would be W itself, hence an abelian variety containing zero and generating \(J_X\). Then \(J_X\) would have dimension two, contrary to nonhyperellipticity. A one-dimensional translate would, by the quasi-finite bijection off zero, admit a finite radicial cover mapping nontrivially to X. Its normalization has genus one (purely inseparable extensions of the elliptic function field preserve genus), so it cannot map nontrivially to a genus\textgreater1 curve; after factoring off inseparability, the separable part contradicts Riemann--Hurwitz.

An exceptional pair has zero abelian Kummer difference after a finite constant extension. Therefore its nonzero difference class lies in \(J_X(\bar k)\{\Sigma'\}\), the torsion supported on the complementary primes. Here \({\Sigma}'\) is finite. The opened proof applies Boxall's finite-prime torsion intersection theorem to W and concludes that its intersection with that subgroup is finite. Injectivity of \({\delta}\) away from the diagonal then gives finitely many exceptional geometric pairs on X. For the prime-to-p choice the complementary prime set is exactly \{p\}, with no missing restriction to odd p, p-rank, ordinarity, or nonhyperellipticity of the original curve. The cover reduction is what handles genus two and hyperelliptic original curves.

We have checked the derivation and the hypotheses of the stated ST theorem. We have not independently reproved Boxall's theorem or the nonhyperelliptic-cover existence result cited within ST's proof.


\end{proof}

\begin{claim}{A.CR4}{Proved}{A rational point section is self-normalizing}

Let $C/L$ be smooth proper geometrically connected of genus at least
two over a finite field, and let $\Sigma$ contain a prime
$r\ne\operatorname{char}L$. If $D\subset\Pi_C^\Sigma$
is a continuous section subgroup mapping isomorphically onto $G_L$,
then $N_{\Pi_C^\Sigma}(D)=D$.

\end{claim}

\begin{proof}

The exact ST2009 Lemma 1.7 hypotheses are: \(Z{\subset}{\Pi}_U\) is closed, pr(Z) is open in \(G_k\), and \(pr|_Z\) is injective. Its conclusion is that pr embeds the commensurator in \(G_k\); commensurator, normalizer, and centralizer coincide and contain Z with finite index. If Z is the section of an L-rational point and we work in \({\Pi}_{U_L}\), its projection is the whole \(G_L\). Therefore the conclusion specializes to

\[N_{\Pi_{U_L}}(Z)=Z_{\Pi_{U_L}}(Z)=Z.\]

This does not assert self-normalization of every closed-point subgroup in the original \({\Pi}_U\) without enlarging constants. It is exactly sufficient for AS-AG.40.

For clarity, the key proof of the lemma can be re-derived. A geometric commensurator \({\sigma}\) centralizes the open intersection of Z and \({\sigma}Z{\sigma}^{-1}\), because both project injectively to the same abelian arithmetic group. If \({\sigma}{\ne}1\), choose an open characteristic \(N{\trianglelefteq}{\Delta}_U\) not containing \({\sigma}\). The open subgroup \(\bar H=\langle N,\sigma\rangle\) has a nonzero image of \({\sigma}\) in its abelianization, detected by the cyclic quotient \(\bar H/N\). The open part of Z centralizing \({\sigma}\) stabilizes \(\bar H\), so this abelian image is fixed by the Frobenius action on the corresponding genuine cover. The abelian geometric group is a product of prime-adic Tate modules of the generalized Jacobian (proper Jacobian for proper covers); none has a nonzero Frobenius-fixed element because its eigenvalues have nonzero Weil weights, and its Tate module is torsion-free. This contradiction proves that the geometric part of the commensurator is trivial. Its injection into the procyclic arithmetic group gives the rest.

The opened source used in the first argument is
\cite[Lemma 1.7]{ST2009}; it has no cofiniteness assumption on \(\Sigma\).


\end{proof}

\begin{claim}{A.CR5}{Proved}{Point recognition in the whole cuspidal quotient}

Under the hypotheses of AS-AG.40, a cusp-preserving isomorphism
$\beta:G_P\simeq G_Q$ over $\mathrm{id}_{\Pi_C}$ carries every
noncuspidal point section, after its finite residue-field extension,
to an actual noncuspidal point section. Outside $E_C\cup\{P,Q\}$ the
target point is the same point of $C$. This conclusion holds in the
whole quotient, and not merely at a finite level.

\end{claim}

\begin{proof}

Suppose \(\beta:G_P\xrightarrow\sim G_Q\) is the isomorphism in AS-AG.40, over \(\operatorname{id}_{{\Pi}_C}\), preserving the cusp decomposition conjugacy class. For a finite extension L/k, restrict both groups to the preimage of \(G_L\); \({\beta}\) remains an isomorphism over \(\operatorname{id}_{{\Pi}_{C_L}}\) and carries the base-changed cusps. Choose \(R{\in}(C-P)(L)\) and a decomposition subgroup \(D_R^P\). Let \(S={\beta}(D_R^P)\), a closed subgroup mapping isomorphically onto \(G_L\).

Let \(\mathcal H\) be the directed set of all open subgroups \(H{\le}(G_Q)_L\) containing S, ordered by reverse inclusion. All these subgroups have full arithmetic projection, so their connected covers \(X_H\) of \(C_L-Q\) are geometrically connected over L. The source cover for \({\beta}^{-1}(H)\) has a rational point above R because it contains \(D_R^P\). By Statement \statementref{A.CR2}'s point-count equality, \(X_H(L)\) is a nonempty finite set.

For \(H'{\subset}H\) the covering map gives a map \(X_{H'}(L){\to}X_H(L)\). The inverse limit of these sets is nonempty even though no individual transition map is asserted surjective. To prove this, take their product, a compact space. Compatibility along each transition defines a closed condition. Any finite set of compatibility conditions is satisfied by choosing a point on the cover for the intersection of the finitely many subgroups involved and projecting it to the relevant covers. Thus the closed conditions have the finite-intersection property, and compactness supplies a compatible family \((x_H)\).

The component at \(H=(G_Q)_L\) is a point \(T{\in}(C-Q)(L)\). Choose a geometric lift \(\widetilde T\) over T in the universal pro-cover with every projection equal to \(x_H\). This is possible by the same compactness argument in the profinite fibre above T. Its decomposition subgroup \(D_{\widetilde T}\) maps isomorphically to \(G_L\): all these covers are étale at T and the residue field upstairs is \(\bar L\). For every H, rationality of \(x_H\) says that \(D_{\widetilde T}{\cap}H\) surjects onto \(G_L\). Since \(D_{\widetilde T}\) itself projects injectively, this implies \(D_{\widetilde T}{\subset}H\). Therefore

\[
D_{\widetilde T}\subset\bigcap_{H\supset S}H=S.
\]

Both sides project isomorphically onto \(G_L\), so they are equal. Thus \({\beta}\) carries a genuine noncuspidal point decomposition subgroup to the decomposition subgroup of a genuine noncuspidal point of the target, in the whole profinite quotient.

If R lies outside \(E_C\) and P,Q, its proper image \(D_R\) in \({\Pi}_{C_L}\) equals the proper image \(D_T\). The finite-exception uniqueness from Statement \statementref{A.CR3} forces T=R. This works for every closed R of C outside the one fixed finite set, upon taking L=k(R). It does not assume an a priori bijection of rational point sets, and does not identify one finite quotient with the inverse limit.


\end{proof}

\begin{claim}{A.CR6}{Proved}{The complete conditional cuspidal criterion}

Let $k$ be finite of characteristic $p$, $C/k$ smooth proper
geometrically connected of genus at least two, $P,Q\in C(k)$, and
$r\ne p$. Let $G_P,G_Q$ be the maximal cuspidally pro-$r$ quotients
of the arithmetic geometrically prime-to-$p$ groups of $C-P,C-Q$,
relative to their maps to the full proper group $\Pi_C$. If a
continuous isomorphism $\beta:G_P\to G_Q$ over
$\mathrm{id}_{\Pi_C}$ carries the cusp decomposition subgroup to
a conjugate of the cusp decomposition subgroup, then $P=Q$.

\end{claim}

\begin{proof}

For clarity about the orientation-class input, the precise finite-coefficient
facts are as follows. For a smooth proper connected curve over an
algebraically closed field and \(r\) invertible in that field,
\(H^2_{\et}(C,\mu_{r^a})\simeq\mathbf Z/r^a\) via degree, with a point
class sent to one; see \cite[Proposition 14.2 and Theorem 24.1(a)]{MilneEtale}.
The pro-\(\Sigma\) group comparison when \(r\in\Sigma\) can be verified
directly from \cite[Theorem 14.9 and Remark 14.10]{MilneEtale}.
In the tower of all connected geometric finite \(\Sigma\)-covers,
every \(H^1(-,\mathbf Z/r^a)\) class is killed by its corresponding
cyclic torsor, and every \(H^2(-,\mu_{r^a})\) class is killed by an
étale cover of degree divisible by \(r^a\). Such covers exist since
the Jacobian has nonzero \(r\)-adic Tate module. The Galois closures of
the composite covers remain \(\Sigma\)-covers, by the normal-core
argument. Cohomology above degree two vanishes. The spectral sequence
therefore identifies group cohomology with curve cohomology for these
coefficients. The same comparison over the finite ground field follows
by comparing the two Hochschild--Serre spectral sequences.
Inverse limits give the required \(\mathbf Z_r(1)\) statement.
The tame surface presentation, including its canonical inertia words,
is \cite[Expos\'{e} XIII, Corollary 2.12]{SGA1}; its hypotheses are a
smooth proper connected curve over a separably closed field and a
finite set of distinct punctures, with the full prime-to-characteristic
quotient taken. Passing to any \(\Sigma\) containing \(r\) preserves the
pro-\(r\) test used here.

Let k be finite of characteristic p, C/k smooth proper geometrically connected with \(g{\ge}2\), \(P,Q{\in}C(k)\), \(r{\ne}p\). Let \(G_P,G_Q\) be the maximal cuspidally pro-r quotients relative to the full proper geometrically prime-to-p arithmetic group \({\Pi}_C\). Suppose \({\beta}:G_P{\to}G_Q\) is a continuous isomorphism over \(\operatorname{id}_{{\Pi}_C}\) carrying the cusp decomposition group to a conjugate of the cusp decomposition group. We prove P=Q.

\textbf{Principal p-primary difference.}

Projection of the cuspidal decomposition groups gives conjugate point decomposition groups at P,Q in \({\Pi}_C\). A conjugator can be made geometric: multiply it by an element of the target decomposition subgroup with the same arithmetic projection. After passing to the geometric abelianization, the Kummer difference \([P-Q]\) has zero class in \(H^1(G_k,T_\ell\) \(J_C)\) for every \(\ell{\ne}p\). Since \(J_C(k)\) is finite, its ell-adic completion is its ell-primary subgroup and the Kummer map is injective on that subgroup. Thus \([P-Q]\) has p-power order. There is \(e{\ge}0\) with \(p^e(P-Q)\) principal. A rational generator over \(\bar k\) descends to \(k(C)^{\times}\) by Hilbert 90: its Galois ratios are constants, and killing that scalar cocycle gives f with

\[\operatorname{div}(f)=p^e(P-Q).\]

\textbf{The central extensions and the multiplier.}

Let \(M=\mathbf{Z}_r(1)\). Divide the geometric cuspidal kernel in \(G_P\) by its commutators with the full geometric subgroup. The resulting arithmetic extension \(E_P{\to}{\Pi}_C\) is the maximal geometrically cuspidally central quotient. Its kernel is procyclic, generated by the puncture inertia, because the original cuspidal kernel is normally generated by the conjugate puncture inertia groups. It is exactly M, not a finite quotient: in the geometric punctured surface presentation the puncture word is the product of the g handle commutators. Map the first pair of handles to \(1+e_{12},1+e_{23}\) in the Heisenberg pro-r group \(\operatorname{UT}_3(\mathbf{Z}_r)\) and all other handles to 1. The puncture word maps to \(1+e_{13}\), of infinite r-adic order, and its normal closure is central in that target. This proves that the central inertia generator has infinite r-adic order. The arithmetic action is the cyclotomic action, as on tame inertia. The same construction gives \(E_Q\).

The class of the geometric central extension is the degree-one orientation class in \(H^2(\Delta_C,M)\simeq Z_r\), with the sign fixed by the Kummer residue convention. One equivalent construction is the Kummer first Chern class of \(O_C(P)\); its canonical trivialization on \(C-P\) gives the homomorphism from the punctured group to the extension and maps local inertia with residue +1. Its geometric restriction is the class of a point, hence a generator. This is the familiar punctured-surface orientation extension; no configuration-group injection enters.

Because \({\beta}\) preserves the proper quotient and the geometric subgroup, it descends to an isomorphism \(\bar\beta:E_P\to E_Q\) over \(\operatorname{id}_{{\Pi}_C}\). On M it acts by a unit \(u{\in}\mathbf{Z}_r^{\times}\). Restricting the extension isomorphism to the geometric group gives u times the degree-one class equal to the same degree-one class. Thus u=1.

\textbf{The difference cocycle and its residues.}

On the continuous fibre product \(E_P{\times}_{{\Pi}_C}E_Q\) define

\[\chi(u,v)=\bar\beta(u)v^{-1}\in M.\]

The common arithmetic projection acts on M cyclotomically. Multiplication verifies \(\chi(ee')=\chi(e)+e\cdot\chi(e')\); therefore \({\chi}\) is a continuous crossed homomorphism. Pull back to the arithmetic geometrically prime-to-p group \({\Gamma}\) of \(C-\{P,Q\}\), using its natural maps to both central extensions. On inertia at P the second map is trivial and the first has residue +1; on inertia at Q the first map is trivial and the inverse of the second gives residue \(-1\). Thus \({\chi}\) has residues \((+1,-1)\).

The Kummer class \({\kappa}_r(f)\) has residues \((p^e,-p^e)\). Hence \({\kappa}_r(f)-p^e{\chi}\) has zero restriction to both inertia groups. Since M is trivial under the geometric group, a cocycle vanishing on these inertia groups vanishes on their closed normal closure. The kernel of \({\Gamma}{\to}{\Pi}_C\) is precisely that closure. Thus this difference descends to \(H^1({\Pi}_C,M)\).

Inflation-restriction gives an isomorphism

\[H^1(G_k,M)\xrightarrow\sim H^1(\Pi_C,M),\]

because \(H^1(\Delta_C,M)^{G_k}=0\). Explicitly \(H^1(\Delta_C,M)\) is \(\operatorname{Hom}(T_rJ_C,\mathbf{Z}_r(1))\); the arithmetic Frobenius eigenvalues on its \(\mathbf{Q}_r\)-extension have complex absolute value \(\sqrt{\#k}\), so none is 1, and the \(\mathbf{Z}_r\)-module is torsion-free. Finite-field Kummer theory identifies \(H^1(G_k,M)\) with the r-primary completion of \(k^{\times}\). Therefore there is \(a{\in}k^{\times}\) representing \({\kappa}_r(f)-p^e{\chi}\).

\textbf{Evaluation at all good closed points.}

Take any closed point R outside \(E_C{\cup}\{P,Q\}\), and put L=k(R). Statement \statementref{A.CR5} gives a conjugacy in \((G_Q)_L\) between \({\beta}(D_R^P)\) and \(D_R^Q\). The projection of a conjugator to \({\Pi}_{C_L}\) normalizes \(D_R\). Statement \statementref{A.CR4} says this normalizer is exactly \(D_R\). Multiply the conjugator by the inverse of the element of \(D_R^Q\) with that projection. This leaves its conjugation of the subgroup unchanged and makes it belong to the cuspidal kernel.

In \(E_Q\) its image therefore lies in M. Conjugating a section by an element of M changes its M-valued cocycle by a coboundary. It follows that \({\chi}\) restricts to zero in \(H^1(G_L,M)\) along the genuine section of R. Evaluation of the preceding Kummer identity gives

\[\kappa_{r,L}(f(R)/a)=0.\]

The point R is allowed to have arbitrary degree. The field L is chosen separately for each R, and \({\beta}\) and the extension comparison are restricted along \(G_L{\to}G_k\); no incompatible Frobenius convention is introduced.

\textbf{Infinitely many usable points and contradiction.}

Assume \(P{\ne}Q\). Then f is nonconstant. Over \(\bar k\) it gives a surjective finite map to the projective line. The excluded closed points have finitely many geometric inverse images, and their images under f form a finite set. Because \(r{\ne}p\), \(\mu_{r^\infty}(\bar k)\) is infinite. Choose a nontrivial \({\zeta}\) in that group with \(a{\zeta}\) outside the finite set of excluded values. Every geometric point over \(a{\zeta}\) is defined over a finite field. Choose one and regard it as a closed point R of C. Then R avoids all exclusions, and \(f(R)/a={\zeta}\) belongs to \(k(R)^{\times}\).

The finite-field Kummer map to \(H^1(G_{k(R)},\mathbf{Z}_r(1))\) is injective on its r-primary subgroup. Consequently the nontrivial r-power root \({\zeta}\) has nonzero Kummer class, contrary to the evaluation identity. This proves P=Q.

This last argument does not require choosing all roots of unity in the original k, and a root becoming an r-th power in some finite extension does not make its entire r-adic Kummer class vanish. The completion of a finite cyclic group preserves its r-primary subgroup.


\end{proof}

\begin{claim}{A.CR7}{Proved}{Proper-field support suffices for the criterion}

Let $\Sigma$ be any set of primes containing $r\ne p$.
Use the full proper geometric pro-$\Sigma$ group in the preceding
criterion. It suffices that the noncuspidal point transport induced
by $\beta$ have unchanged labels outside $C(F)$, for one proper
algebraic subfield $k\subset F\subsetneq\bar k$. In particular,
an exceptional image supported over such an $F$ suffices, even if it
is infinite. The annihilating integer for $[P-Q]$ need not be
coprime to $r$.

\end{claim}

\begin{proof}

This strengthening also applies to the \(J_C\)-large prime sets of \cite{ST2017}.

Fix any prime set \({\Sigma}\) containing a prime \(r\ne p\), and use the full proper geometric pro-\({\Sigma}\) group in the definitions of \(G_P,G_Q\). The same-prime lemma, compactification comparison, central extension construction, and self-normalizer argument remain valid. Assume \({\beta}\) is a cusp-preserving isomorphism over \(\operatorname{id}_{{\Pi}_C^{\Sigma}}\). Suppose that there is a proper algebraic subfield \(F{\subsetneq}\bar k\) containing k such that, for every closed R whose geometric points are not F-rational, the point transported by \({\beta}\) has the same label R. Then P=Q.

Only the final Kummer argument changes. Let n be the exact order of \([P-Q]\) in \(J_C(k)\), and choose \(f{\in}k(C)^{\times}\) with \(\operatorname{div}(f)=n(P-Q)\). This order exists because \(J_C(k)\) is finite. The argument for the central extensions uses neither that n is p-primary nor that n is coprime to r: the difference \({\kappa}_r(f)-n{\chi}\) has zero residues, hence is a constant class \({\kappa}_r(a)\), and \({\chi}\) evaluates to zero at every point with unchanged label. Thus \({\kappa}_r(f(R)/a)=0\) for every such R.

The function f/a is not an r-th power in \(\bar k(C)^{\times}\). If \(r\nmid n\) its valuations at P,Q already rule this out. If \(r\mid n\) and \(f/a=b^r\), then \(\operatorname{div}(b)=(n/r)(P-Q)\), contradicting the exact order of \([P-Q]\). The geometric divisor class has the same order as its k-rational class, since \(J_C(k){\to}J_C(\bar k)\) is injective. Equivalently, a rational principal divisor descends by Hilbert 90.

Let D be the smooth compactification of the Kummer cover \(z^r=f/a\). It is geometrically connected, since r is prime and f/a is not a geometric r-th power; \(r{\ne}p\) makes it generically separable. Take a finite L/k containing \({\mu}_r\). Because F is a proper subfield of \(\bar k\), some prime ell has finite exponent in the supernatural degree {[}F:k{]}. To see this directly, the nonzero closed subgroup \(G_F{\subset}\widehat{\mathbf Z}\) has a nonzero ell-adic projection for some ell, and a nonzero closed subgroup of \(\mathbf{Z}_\ell\) is open. Put \(L_m\) equal to the degree-\(\ell^m\) extension of L. Then

\[F_0=F\cap\bigcup_m L_m\]

is a finite extension of k: the degrees outside ell are bounded by {[}L:k{]}, and the ell-part of {[}F:k{]} is bounded. Thus all excluded F-rational points lying in \(C(L_m)\) belong to the fixed finite set \(C(F_0)\). Write \(B=\#C(F_0)+2\), including P,Q.

For every remaining \(R{\in}C(L_m)\), the zero r-adic Kummer class implies that f(R)/a is an r-th power in \(L_m\). Since \({\mu}_r{\subset}L_m\) and f has no zero or pole away from P,Q, the fibre of D above R consists of exactly r \(L_m\)-rational points. Hence

\[\#D(L_m)\ge r\bigl(\#C(L_m)-B\bigr).\]

Writing \(Q_m=\#L_m\) and applying the Weil bounds to the geometrically connected proper curves C and D gives

\[(r-1)(Q_m+1)\le (2g_D+2r g_C)\sqrt{Q_m}+rB.\]

This is impossible as \(m{\to}{\infty}\). Thus P=Q. The auxiliary ell is allowed to equal p or r; the extension is an extension of finite constants and no ramification issue occurs. Only the Kummer degree r must be different from p.

In particular, if \(E_C^{\Sigma}\) is supported over a proper algebraic subfield, the needed same-label assertion follows from compactness recognition and the definition of E. Thus, once the generalized configuration transport is established, almost \({\Sigma}\)-separated in the sense of ST2017 implies \({\Sigma}\)-separated, i.e.~\(E_C^{\Sigma}=\varnothing\). This conclusion uses no finiteness assertion about \(E_C\) itself.

The finite-exception case is a special case: finitely many closed-point residue fields generate one finite F, which is proper in \(\bar k\). In fact the original roots-of-unity proof also works directly with arbitrary n, so no coprimality condition was hidden there.


\end{proof}

\begin{claim}{A.CR8}{Proved}{An independent elementary support-field proof}

For a proper hyperbolic curve $C/k$ over a finite field of
characteristic $p$, put
\[
B_C=\{R\in C(\bar k):\exists S\ne R,
       [R-S]\in J_C(\bar k)[p^\infty]\}.
\]
There is a proper algebraic subfield $F\subsetneq\bar k$ containing
$k$ for which $B_C\subset C(F)$. Hence the conditional criterion
\statementref{A.CR6} admits a proof using
\statementref{A.CR7} and this support field in place of the
exceptional-set finiteness theorem.

\end{claim}

\begin{proof}

We can remove the Boxall finiteness input entirely from an alternative proof of the main point-separation theorem. Define

\[B_C=\{R\in C(\bar k):\exists S\ne R,\ [R-S]\in J_C(\bar k)[p^\infty]\}.\]

We prove that \(B_C\) is supported over one proper algebraic subfield of \(\bar k\). This is a statement about divisor classes, not a replacement of group collisions by a claimed converse.

\textbf{The fibres of the difference map.}

For every nonzero \(t{\in}J_C(\bar k)\), the set

\[\{(R,S)\in C(\bar k)^2:R\ne S,[R-S]=t\}\]

has at most two elements. Indeed, for two such pairs, R+S' is linearly equivalent to R'+S. If these effective degree-two divisors are equal, the nonzero-difference condition forces the ordered pairs to be equal. Otherwise they define a degree-two pencil on C. There is at most one such pencil when \(g{\ge}2\): two distinct degree-two maps would give a birational map of C to a curve of bidegree (2,2) in \(P^1{\times}P^1\), whose arithmetic genus is one, impossible. The degree-two pencil is separable even in characteristic two, because a purely inseparable degree-two map to \(P^1\) would identify C with a Frobenius twist of a rational curve, again impossible in positive genus. Denote its involution by iota. Equality of the two fibre divisors then forces (R',S')=(iota(S),iota(R)). This is the sole possible second pair.

\textbf{A nontrivial arithmetic subgroup fixing all p-primary torsion.}

Let f be the p-rank of \(J_C\). Then \(T_pJ_C\cong Z_p^f\) and its finite quotients identify with the geometric p-power torsion groups. Choose an odd prime \(\ell{\ne}p\) not dividing \(|\operatorname{GL}_f(\mathbf{F}_p)|\), with \(|\operatorname{GL}_0(\mathbf{F}_p)|=1\). The kernel of \(\operatorname{GL}_f(\mathbf{Z}_p){\to}\operatorname{GL}_f(\mathbf{F}_p)\) is pro-p, as its congruence filtration has elementary abelian p-group successive quotients. The ell-adic factor \(H=\mathbf{Z}_\ell{\subset}G_k\) has trivial image in \(\operatorname{GL}_f(\mathbf{F}_p)\), since the latter order is prime to ell; its image in the pro-p congruence kernel is also trivial. Thus H fixes all of \(J_C(\bar k)[p^{\infty}]\).

If \(t=[R-S]\) is a nonzero p-primary difference, H acts on the at-most-two-element fibre over t. A pro-ell group with ell odd cannot act nontrivially on a set of cardinality at most two. Hence H fixes R and S individually. Therefore

\[B_C\subset C(F),\qquad F=\bar k^H\subsetneq\bar k.\]

If f=0, the p-primary torsion group itself is zero, so \(B_C\) is empty because the Abel map of a positive-genus curve is injective.

\textbf{Application to the conditional criterion.}

From the compactness proof, \({\beta}\) carries the point R to an actual target point S over the same finite field. Because \({\beta}\) is over \(\operatorname{id}_{{\Pi}_C}\), their full proper point sections have equal geometric conjugacy class. Passing to the Jacobian Kummer map proves \([R-S]\) is p-primary. This is only the necessary implication from group collision to divisor-class torsion, which has already been proved in the Kummer step of \statementref{A.CR6}. Therefore R outside \(B_C\) must have S=R.

Statement \statementref{A.CR7} now applies with the explicit proper field F above and proves P=Q. This independently establishes the conditional cuspidal criterion without ST2009's exceptional-set finiteness, Boxall's theorem, or an assumption that a p-primary divisor difference produces a group collision. It retains only the exact cusp isomorphism as input, the full same-prime compactification lemma, the compactness recognition, the elementary degree-two-pencil bound, standard curve cohomology, finite-field Kummer theory, and the Weil bound.

Consequently, after the independent verification of AG.41 and AG.42, this is an alternative complete route to the main theorem. The original K1 audit remains useful for checking the earlier proof exactly as written, but it is not an indispensable external premise of this replacement proof.


\end{proof}

\begin{claim}{A.CR9}{Proved}{Exceptional emptiness and all-extension point separation}

For a proper hyperbolic $C/k$ over a finite field of characteristic
$p$ and every prime set $\Sigma$ containing some $r\ne p$,
$E_C^\Sigma=\varnothing$ is equivalent to injectivity of the
geometric pro-$\Sigma$ point-section map after every finite
constant extension of $k$.

\end{claim}

\begin{proof}

For every prime set \({\Sigma}\) containing some prime \(r\ne p\), let \({\Pi}_C^{\Sigma}\) retain the full arithmetic quotient. Assume C is proper hyperbolic over finite k. Then \(E_C^{\Sigma}\) is empty if and only if the point-section map is injective after every finite constant extension.

If distinct L-rational points have geometrically conjugate sections, choose their lifts so that the subgroups are equal. The universal pro-cover for \(C_L\) is the same universal pro-cover with the open arithmetic group \({\Pi}_{C_L}^{\Sigma}\). Equal decomposition groups are commensurate. The lifts remain distinct because their geometric images are distinct. Thus the original exceptional set is nonempty.

Conversely, take two distinct upstairs points with commensurate decomposition groups \(D_1,D_2\). Proposition 1.8(iii), whose proof only needs Lemma 1.7 and no cofiniteness, shows that their geometric images are distinct. The intersection \(D_1{\cap}D_2\) is open in both. Its arithmetic projection is open in \(G_k\) because the \(D_i\) project injectively onto open subgroups. Choose a finite L/k containing both point residue fields and with \(G_L\) contained in that projection. Since each \(D_i\) projects injectively, their restrictions over \(G_L\) are precisely the same inverse image of \(G_L\) in \(D_1{\cap}D_2\). They give equal sections of distinct L-points. This is a collision in the whole \({\Pi}_{C_L}^{\Sigma}\).

ST2017 Proposition 2.11(i), read in full, states \(E_{C'}=f^{-1}(E_C)\) for any étale \({\Sigma}\)-covering f, including constant extensions. This is compatible with the equivalence. Its ramified statement is an inclusion with the ramification locus, so must not be promoted to equality. Proposition 2.12, under \(J_C\)-largeness, only confines exceptional points to a proper algebraic subfield, finite of degree at most two over the \({\Sigma}'\)-torsion field (degree one for nonhyperelliptic C); it does not exhibit a nonempty exceptional set or contradict emptiness for the full prime-to-p set.


\end{proof}

\begin{claim}{A.CR10}{Proved}{Large prime sets and the finite-or-infinite dichotomy}

Assume the cuspidal transport construction from equal point sections
has been established for a prime set $\Sigma$ containing some
$r\ne p$, allowing $p\in\Sigma$. Then
$E_C^\Sigma=\varnothing$ whenever $\Sigma$ is $J_C$-large in the
sense of \cite{ST2017}; this includes every cofinite $\Sigma$.
More generally, almost $\Sigma$-separation implies
$E_C^\Sigma=\varnothing$, and a finite $E_C^\Sigma$ is necessarily
empty. Thus the exceptional image is either empty or infinite.
This dichotomy alone does not classify every prime set.

\end{claim}

\begin{proof}

For any prime set \({\Sigma}\) containing some \(r\ne p\), put

\[T_{\Sigma'}J_C=\prod_{s\in\Sigma'}T_sJ_C,\qquad H=\ker(G_k\to\operatorname{Aut}(T_{\Sigma'}J_C)).\]

Here the p-adic Tate module means the étale p-adic module, of rank equal to the p-rank. Under the ST2017 definition, \({\Sigma}\) is \(J_C\)-large exactly when H is nontrivial. Every \({\Sigma}'\)-primary geometric torsion point is fixed by H. The fibre argument of the difference-map argument in \statementref{A.CR8} shows that 2H fixes each ordered pair realizing a nonzero \({\Sigma}'\)-primary difference. Since multiplication by two is injective on the closed subgroup \(H{\subset}\widehat{\mathbf Z}\), 2H is still nonzero. Thus

\[B_C^\Sigma\subset C(\bar k^{2H}),\qquad \bar k^{2H}\ne\bar k.\]

This recovers and makes explicit the support mechanism of ST2017 Proposition 2.12: for a nonhyperelliptic curve the fibre is a singleton, so H itself suffices; for a hyperelliptic curve one may need the index-at-most-two subgroup 2H.

The Kummer assertion used here includes the characteristic prime when
$p\in\Sigma$. Over a finite field $L$, multiplication by $p^n$ is
surjective on $J_C(\bar L)$, with kernel $J_C[p^n](\bar L)$.
It gives an injective map
\[
 J_C(L)/p^nJ_C(L)\longrightarrow
 H^1(G_L,J_C[p^n](\bar L)).
\]
These are the same classes as those read by the abelian geometric
pro-$p$ fundamental group. Indeed, factor $[p^n]$ through the
quotient by the connected component $J_C[p^n]^0$: the first
isogeny is radicial and bijective on geometric points, and the
remaining isogeny is finite \emph{\'etale}, with geometric kernel
$J_C[p^n](\bar L)$. Its lifting cocycle therefore equals the
multiplication-$p^n$ Kummer cocycle. The Abel--Jacobi map identifies
the abelian curve group with the Jacobian group at this prime as
well, by \cite[Proposition 9.1]{MilneJac}. Passing to the inverse
limit is legitimate for these finite groups and detects the entire
$p$-primary part of the finite group $J_C(L)$. Thus equality of
point sections kills the components at every prime in $\Sigma$,
including $p$ when it is retained; the remaining difference has
precisely complementary-prime torsion. No ordinary-Jacobian
hypothesis is required.

In the generalized criterion, compactness carries each R to S with \([R-S]\) \({\Sigma}'\)-primary by Kummer theory at all primes in \({\Sigma}\). Hence outside \(B_C^{\Sigma}\) the labels agree. Statement \statementref{A.CR7} then proves the generalized conditional criterion for every \(J_C\)-large \({\Sigma}\). If the configuration transport is verified for every such \({\Sigma}\), it follows that \(E_C^{\Sigma}=\varnothing\) for every \(J_C\)-large \({\Sigma}\) and every proper hyperbolic C over a finite field. Every cofinite \({\Sigma}\) is \(J_C\)-large: choose an auxiliary odd ell outside the finite set \({\Sigma}'\) and outside the prime divisors of all finite groups \(\operatorname{GL}_{\operatorname{rank}T_s}(\mathbf F_s)\), \(s{\in}{\Sigma}'\); then \(\mathbf{Z}_\ell\) acts trivially on the complementary Tate module.

Even without \(J_C\)-largeness, finiteness of \(E_C^{\Sigma}\) supplies finite-field support and Statement \statementref{A.CR7} applies. Therefore, again under the generalized transport theorem, a finite exceptional set must be empty. Every \(E_C^{\Sigma}\) is consequently either empty or infinite. This is a dichotomy, not a classification of which alternative occurs for every arbitrary \({\Sigma}\). The singleton theorem in Statement \statementref{A.CR13} supplies explicit infinite cases; the \(J_C\)-large theorem supplies empty cases. We have not proved that these exhaust all possible prime sets.


\end{proof}

\begin{claim}{A.CR11}{Proved}{Genus-one negative control}

The elliptic curve $E/\mathbf F_2:y^2+xy=x^3+1$ has distinct
rational points $O$ and $T=(0,1)$ whose sections collide in the
entire geometrically prime-to-two arithmetic group. Its exceptional
image is all of $E^{\cl}$. The finite-exception hypothesis used for
proper hyperbolic curves fails here.

\end{claim}

\begin{proof}

Take \(E/\mathbf{F}_2\) given by \(y^2+xy=x^3+1\), with origin O at infinity. Its affine \(\mathbf{F}_2\)-points are (0,1),(1,0),(1,1), so it has four rational points. It is nonsingular: a potential affine singularity must have x=0 from \(\mathbf{F}_y=x\), but at (0,1) the derivative \(\mathbf{F}_x=y+x^2\) is 1; the standard projective infinity point is also nonsingular. The inverse of (x,y) is (x,y+x). Hence T=(0,1) is a nonzero point of exact order two.

The entire geometric fundamental group of an elliptic curve is abelian. Its prime-to-two quotient is the product of the odd Tate modules. Using O as splitting, the class of the point section of T is exactly its Kummer class in that product. The two-primary point T is divisible by every odd integer in the finite group \(E(\mathbf{F}_2)\); consequently every odd-adic Kummer component vanishes. Thus O,T collide in the entire geometrically prime-to-two group, not just in an abelianized approximation of a larger nonabelian group.

For every \(x\in E(\bar F_2)\), the same argument over a finite field containing x gives a collision between x and x+T. These are distinct. Thus the exceptional image is all of \(E^{cl}\); it is infinite. The exact failed AS-AG.40 hypothesis is the finite-exception statement in Step 3. ST2009 Proposition 1.8 assumes hyperbolicity of U; the proper genus-one curve is not hyperbolic. The rest of the Kummer argument does not repair that missing input.


\end{proof}

\begin{claim}{A.CR12}{Proved}{The abelianized negative control in genus two}

The smooth proper genus-two curve
$C/\mathbf F_2:y^2+x(x+1)y=x^5+1$ has distinct rational points
$P=(0,1)$ and $Q=(1,0)$ with $[P-Q]$ of exact order two. They
collide in the entire arithmetic extension by the geometric
abelian prime-to-two quotient. Both the cusp-residue and the
compactified-cover cohomology arguments cease to apply to that
abelianized replacement.

\end{claim}

\begin{proof}

Take \(C/\mathbf{F}_2\) given by

\[y^2+x(x+1)y=x^5+1.\]

At finite ramification points x=0,1 there are respectively P=(0,1),Q=(1,0), and \(\mathbf{F}_x=y+x^4\) is 1 at both. The infinity chart \(u=1/x,v=y/x^3\) is

\[v^2+(u+u^2)v=u+u^6,\]

whose derivative in u is 1 at (0,0). This is a smooth genus-two hyperelliptic curve, with three simple Artin--Schreier branch points \((0,1,{\infty})\). For an elementary genus check, put z=y/{[}x(x+1){]} in the function field. The right-hand Artin--Schreier function is \((x^5+1)/(x^2(x+1)^2)\). At 0, subtract the Artin--Schreier coboundary \((1/x)^2+1/x\); the remaining pole has order one. At 1 the original pole has order one, and at infinity it has order one. The degree-two separable map therefore has different contribution two at each of three points, giving \(2g-2=-4+6=2\).

The rational function x/(x+1) has divisor \(2P-2Q\). The class \([P-Q]\) is nonzero, because a function with divisor \(P-Q\) would give a degree-one map to the projective line, impossible in positive genus. Hence its exact order is two. Both point sections become equal up to geometric conjugation in the entire extension by \(\Delta_C^{\rm ab,(2')}\), by finite-field Kummer theory. The program gives \(N_1=N_2=3\) and \(\#J_C(\mathbf{F}_2)=4\), consistently.

The full AS-AG.40 input does not survive replacing all groups by their geometric abelianizations. At a once-punctured curve the sole puncture inertia is zero in the geometric abelianization, since the puncture word is a product of commutators. Thus there is no central cusp kernel \(\mathbf{Z}_r(1)\) in that abelianized punctured group and no cocycle with residue +1 at that cusp. Moreover, even the cohomology-retention statement needed for point recognition fails for general open subgroups of the proper abelian quotient: an index-r geometric subgroup has r-adic abelian rank 2g, while its actual degree-r étale curve cover has genus \(1+r(g-1)\), hence r-adic rank \(2+2r(g-1)>2g\). Therefore the relevant failure occurs before using finite exceptions or Kummer evaluations; AG.39.15 must be applied to the full punctured pro-prime-to-p group and its actual cuspidal kernel, not to an abelianized replacement.


\end{proof}

\begin{claim}{A.CR13}{Proved}{Empty and singleton prime sets have infinitely many exceptions}

For every proper hyperbolic curve $C/k$ over a finite field and every
prime $r$, including $r=\operatorname{char}k$, the exceptional
image $E_C^{\{r\}}$ is infinite. Here the full arithmetic quotient
$G_k$ is retained, and for $r=\operatorname{char}k$ we use the
same definition by commensurable point decomposition groups.
For $\Sigma=\varnothing$, the group is $\Pi_C^\varnothing=G_k$;
all rational point sections are equal and every closed point is
exceptional. Thus retaining at most one geometric prime does not
separate points over every finite constant extension.

\end{claim}

\begin{proof}

\textbf{Elementary automorphism lemma.} Let D be a finitely generated pro-r group, with Frattini quotient \(D/\Phi(D)\cong F_r^d\). Then the kernel of \(\operatorname{Aut}(D){\to}\operatorname{GL}_d(\mathbf{F}_r)\) is pro-r.

Here is a proof avoiding an unopened automorphism theorem. For a finite r-group H, every maximal subgroup has index r, so a tuple whose images generate \(H/{\Phi}(H)\) generates H. Fix one basis of \(H/{\Phi}(H)\), and consider the set of all ordered lifts of that basis to H. It has cardinality \(|{\Phi}(H)|^d\), an r-power. The kernel of \(\operatorname{Aut}(H){\to}\operatorname{Aut}(H/{\Phi}(H))\) acts freely on this set: fixing a generating tuple fixes H pointwise. Therefore the kernel's order divides this r-power, so it is an r-group. A finitely generated pro-r group has a cofinal system of open characteristic subgroups (intersect all subgroups of bounded index). Its finite quotients are r-groups; the image of \({\Phi}(D)\) in each is the finite quotient's Frattini subgroup. Passing to the inverse limit of the preceding finite kernels proves the assertion.

Now assume C/k proper hyperbolic has a rational point O, write \({\Delta}_r\) for its geometric pro-r group, and use \(s_O\) to express \({\Pi}_C^{\{r\}}={\Delta}_r{\rtimes}G_k\). Its Frattini rank is \(d=2g\) when \(r\ne p\), and \(d=f_C\), the geometric p-rank of the Jacobian, when \(r=p\). The latter equality follows from the Artin--Schreier comparison and the all-prime Abel--Jacobi comparison in \cite[Proposition 9.1 and Lemma 9.4]{MilneJac}. In particular it is finite. For a pro-r group, finitely many lifts of a basis of its Frattini quotient generate the group: otherwise their closed generated subgroup lies in a maximal open subgroup of index r, contradicting that they generate the quotient. Thus the automorphism lemma applies at the characteristic prime as well. If \(f_C=0\), the pro-p group itself is trivial. Choose a prime q different from r and not dividing \(|\operatorname{GL}_d(\mathbf{F}_r)|\), with \(\operatorname{GL}_0\) interpreted as the trivial group. The q-adic direct factor \(Z_q\subset\widehat Z=G_k\) has trivial image under \(G_k{\to}\operatorname{Aut}({\Delta}_r)\): its image in \(\operatorname{GL}_d(\mathbf{F}_r)\) is a q-group, hence trivial, and its remaining image in the pro-r kernel is both pro-q and pro-r, hence trivial. Consequently \(s_O(\mathbf{Z}_q)\) is central in the entire arithmetic group.

For every closed point x of degree divisible by q, choose \(h{\in}\mathbf{Z}_q\) with q-adic coordinate 1. This h acts nontrivially on a geometric point over x, because its image modulo deg(x) is nonzero. Let \(z=s_O(h)\), and choose any lift \(\tilde{x}\) in the universal pro-cover. Then \(z\tilde{x}\) differs from \(\tilde{x}\) already on the geometric curve, whereas

\[D_{z\tilde{x}}=zD_{\tilde{x}}z^{-1}=D_{\tilde{x}}.\]

Thus x is exceptional. There are infinitely many such closed points: for sufficiently large j the Weil bound implies \(\#C(k_{q^j})>\#C(k_{q^{j-1}})\), so there is a point of exact degree \(q^j\). The equality of the displayed decomposition groups, restricted to a finite field defining both geometric points, gives equal whole section subgroups for two distinct rational points. No finite-level approximation is involved.

If C(k) is empty, first take a finite constant extension with a point. Equality of two decomposition groups in this open arithmetic subgroup makes their decomposition groups in the original arithmetic group commensurable: each original group contains that restricted group with finite index. The geometric points are still distinct. Infinitely many closed points over the extension map to infinitely many closed points over k because the extension is finite. This proves descent of infiniteness directly, without invoking ST2009's additional prime-set hypotheses. Therefore \(E_C^{\{r\}}\) is infinite for every proper hyperbolic C over every finite field and every prime r.

When \(\Sigma=\varnothing\), the geometric quotient is trivial. A geometric point over a closed point of degree n has decomposition group \(G_{k_n}\subset G_k\). Any two such open subgroups of the procyclic group are commensurable. Since the curve has infinitely many geometric points, every closed point is exceptional, and all rational points over any fixed finite extension give the same identity section of that extension's whole arithmetic group.

This proof gives actual all-level collisions after finite extensions. It also shows exactly why ST2009's cofiniteness hypothesis cannot be omitted from AS-AG.40's finite-exception argument. We make no inference that arbitrary finite sets with more than one prime have virtually pro-\({\Sigma}\) automorphism groups; that would require a separate theorem and is not supplied here.


\end{proof}

\begin{claim}{A.CR14}{Proved}{A completely specified pro-three collision}

Let $L=\mathbf F_2[t]/(t^7+t+1)$ and $C/L:y^2+y=x^5$.
The distinct $L$-points
$P=(t,t^6)$ and $Q=(t^2,t^6+t^5)$ have conjugate sections in
the entire arithmetic extension by the geometric pro-three group.
The assertion is not merely a finite-quotient collision.

\end{claim}

\begin{proof}

Let \(L=F_2[t]/(t^7+t+1)\), and take \(C/L:y^2+y=x^5\) with \(O={\infty}\). It is smooth of genus two. The program verifies the field identity for every nonzero element of the 128-element ring, so irreducibility of the chosen modulus is checked rather than assumed. Put

\[P=(t,t^6),\qquad Q=(t^2,t^6+t^5).\]

They satisfy the equation and Q is the \(\mathbf{F}_2\)-Frobenius image of P. The orbit has seven distinct points; the code gives their complete coordinates. Work in the geometrically pro-three quotient over \(\mathbf{F}_2\). We have \(|GL_4(F_3)|=24\,261\,120\), which is not divisible by seven. The preceding argument gives a central image of \(\mathbf{Z}_7{\subset}G_{\mathbf{F}_2}\). Its element with 7-adic coordinate 1 has image 1 in \(\operatorname{Gal}(L/\mathbf{F}_2)\), so sends a lift of P to a lift of Q while leaving the decomposition group unchanged. Intersect that common decomposition subgroup with the preimage of \(G_L\). This gives equal section subgroups for the two L-rational points in the entire geometrically pro-three arithmetic group over L.

The program finds \(\#C(L)=129\), of which 126 points have exact degree seven over \(\mathbf{F}_2\), making 18 exceptional degree-seven closed points. Infinitely many other degrees are handled by the preceding theorem; no bounded computation is being used to prove infiniteness.


\end{proof}

\begin{claim}{A.CR15}{Open}{The remaining arbitrary-prime-set classification}

The assertions above prove emptiness for $J_C$-large sets (once
the stated transport theorem is inserted), infiniteness for
singletons, and a finite-or-infinite dichotomy. They do not decide
the alternative for every arbitrary nonempty prime set that is
neither a singleton nor $J_C$-large. That additional classification
asked in A12 remains unproved in this audit.

\end{claim}

\paragraph{State of the unresolved assertion.}

The remaining obstruction is exact: the Kummer comparison confines a
possible exceptional pair to a fibre over complementary-prime torsion,
but when the complementary Tate representation is faithful its kernel
gives no proper field supporting those pairs. The single-prime central
arithmetic subgroup proof cannot be applied to a multi-prime completion:
its automorphism group need not have a pro-\(\Sigma\) kernel on a finite
Frattini quotient. Neither assertion is silently assumed. The next
approach would be to classify the kernel of the full outer Frobenius
action for these prime sets, or to establish sparse support of the
possible point-pair differences along a constant-field tower.

